\clearpage
\begin{column}{神奈川県央地域でのデマンド型交通の対象地区の選定}
{\small 永井鴻平}
    
　神奈川県央に位置する海老名市, 綾瀬市, 大和市, 座間市の4市を対象に, オンデマンド交通の対象ゾーン選定の検討を行った.

　対象ゾーン内の利用者は, ハブから出発した車両にピックアップされ, 一旦ハブへ集約, その後目的地へ送られることを想定する.
ハブは, 4市のおよそ中央に位置するさがみの駅南口とした.

　対象の需要は, 4市内を相互発着する朝7時から10時台のトリップとし, 鉄道によるものを除いた.
対象のトリップ数は1,723件であった.
選定するゾーンの単位は東京都市圏パーソントリップ調査の地区コードとし, 4市内のゾーン数は合計で319であった.

　ゾーンの選定は, 目的関数
\begin{equation}
    \sum_{i \in I} w_{i} \~{q}_i x_{i} - \nu k
\end{equation}
を最大化する問題として定式化した.
決定変数はゾーン$i$が選定されるか否かを表す二値変数$x_i$と, サービスに供される車両の台数$k$である.
ここで, $I$はゾーン全体の集合, $w_i$はゾーン$i$の重み, $\~{q}_i$はゾーン$i$の有効需要, $\nu$は車両1台あたりのコストである.

　重み$w_i$は, 自身での運転や徒歩での移動が困難な利用者を重視したサービス設計のため, ゾーン$i$の対象トリップのうち, 60歳以上または外出に身体的困難を伴う利用者によるトリップの割合$r_i$を用いて,
\begin{equation}
    w_i = 1 + \theta r_i
\end{equation}
と定義した. $\theta$は外出が困難な利用者をどの程度重視するかを表すパラメータである.
$\theta = 1$, すなわちゾーン$i$の対象トリップが全て外出困難者によるとき, $w_i = 2$となり, そのゾーンの重みが通常の2倍になる.

　有効需要$\~{q}_i$は, 対象トリップのうち, サービスの初期利用者として10\%程度が期待できると仮定し,
\begin{equation}
    \~{q}_i = 0.1 q_i
\end{equation}
と定義した. $q_i$はゾーン$i$を発着する対象トリップ数である.

　車両1台あたりのコスト$\nu$は, サービス対象とする需要とのバランスを考慮し, 1台あたりのコストが需要20件と対応するよう,
$\nu = 20$
とした.

　決定変数の整数制約条件として, 
\begin{align}
    x_i &\in \{0, 1\} \\
    k &\in \mathbb{Z}, 1 \leq k \leq 30
\end{align}
を課した. 
ゾーン$i$が選定されるとき$x_i = 1$, 選定されないとき$x_i = 0$となる.
車両の台数の上限は30台とした.

　また, 容量制約
\begin{equation}
    \sum_{i \in I} \~{q}_i s_i x_i \leq \eta H k,
\end{equation}
およびサービス対象とする需要の最低カバー率制約
\begin{equation}
    \sum_{i \in I} \~{q}_i x_i \geq \gamma \sum_{i \in I} \~{q}_i
\end{equation}
を課した.
ここで, $s_i$はハブからゾーン$i$を往復する時の所要時間であり, 平均走行速度$18 \, \mathrm{km/h}$, 乗り降りや車両の折り返しにかかるバッファ時間を8分と仮定して,
\begin{equation}
    s_i = \varepsilon \left( \frac{2d_i}{18} \times 60 + 8 \right) \, [\text{分}]
\end{equation}
と定義した. $d_i$はハブからゾーン$i$までの走行距離である.
走行距離を求める際, ゾーンの代表点はゾーン内の対象トリップ発着点の重心とした.
$\varepsilon$は, 需要が集中したゾーンでの相乗り効果を考慮するための係数であり, ゾーン$i$の有効需要の密度$p_i\, [\mathrm{件/km^2}]$を用いて,
\begin{equation}
    \varepsilon = 1 - 0.1565 \sqrt{p_i}
\end{equation}
と定義した.
選定対象のゾーンのうち, $p_i$最大のゾーンで$\varepsilon \approx 0.3$となるよう設定した.
$\eta$は車両の稼働率, $H$はサービス運用時間, $\gamma$は需要の最低カバー率である.
需要の対象トリップの抽出時間に合わせ, $H = 240\, [\mathrm{分}]$とした.
また, $\eta = 0.85$, $\gamma = 0.5$とした.

　この問題は, ゾーンの選定と車両台数の決定を同時に行う組合せ最適化問題であり, 整数計画問題として解くことができる.
対象の需要$\~{q}_i$の分布と, $s_i/\varepsilon$の分布, および$d_i$の算出に用いた道路ネットワークを図\ref{col4-5:fig:precondition}に示す.

\begin{center}
    \caption{figure}{需要の分布}
    \label{col4-5:fig:precondition}
\end{center}

　$\theta = 10$としたときのゾーンの選定結果を図\ref{col4-5:fig:result1}に示す.
このとき, $k = 15 \,[\mathrm{台}]$であった.

\begin{center}
    \caption{figure}{ゾーンの選定結果}
    \label{col4-5:fig:result1}
\end{center}

\end{column}
\clearpage